\documentclass{book}

\usepackage{graphicx,epsfig}
\usepackage{hyperref}
\usepackage{amsmath,amssymb}
\usepackage{enumitem}
\usepackage[toc,page]{appendix}
\usepackage{xcolor}
\renewcommand{\baselinestretch}{1}

\setlength{\textheight}{9in}
\setlength{\textwidth}{6.5in}
\setlength{\headheight}{0in}
\setlength{\headsep}{0in}
\setlength{\topmargin}{0in}
\setlength{\oddsidemargin}{0in}
\setlength{\evensidemargin}{0in}
\setlength{\parindent}{.3in}
\pagestyle{plain}

% Helper macro for image placeholders
\newcommand{\imgplaceholder}[2]{%
  \begin{center}%
    \fbox{\parbox{0.72\textwidth}{\centering\vspace{1.4cm}%
      {\large\textbf{[Image Placeholder]}}\\[0.35cm]%
      \textbf{#1}\\[0.2cm]%
      \textit{\small #2}%
      \vspace{1.4cm}}}%
  \end{center}%
}

\begin{document}
\bibliographystyle{plain}

\large

\thispagestyle{empty}
\centerline{\bf A REPORT}
\vspace*{0.3cm}
\centerline{\bf ON}
\vspace*{0.3cm}
\centerline{\bf DEVELOPMENT OF AN AGENTIC SUPPLY CHAIN}
\centerline{\bf INTELLIGENCE PLATFORM USING GILHARI, ORMCP,}
\centerline{\bf LANGGRAPH, AND AZURE CLOUD}
\vspace*{2cm}

\centerline{BY}
\vspace*{1cm}

\begin{center}
	\begin{tabular}{lll}
		Name(s) of the students & \hspace*{2cm} ID No(s). & \hspace*{2cm} Discipline(s) \\
		Sudhanva B H &\hspace*{2cm} 2024A7PS0085H &\hspace*{2cm} B.E. Computer Science \\
	\end{tabular}
\end{center}

\vspace*{2cm}
\begin{center}
	{\bf Prepared in partial fulfillment of the \\
	Practice School - I} \\
\end{center}

\vspace*{0.5cm}

\centerline{AT}

\vspace*{1cm}
        \centerline{\bf Software Tree, California, USA}
	\vspace{0.2cm}
	\centerline{\bf A Practice School - I}
	\vspace{0.2cm}
	\centerline{\bf Station of}
	\vspace{0.5cm}
\begin{figure}[ht]
\epsfxsize=1.0in
\centerline{\epsffile{logo.eps}}
\end{figure}

	\centerline{\bf BIRLA INSTITUTE OF TECHNOLOGY \& SCIENCE, PILANI}
	\vspace*{0.5cm}
	\centerline{\bf July, 2026}
	\newpage

\thispagestyle{empty}

	\centerline{\bf BIRLA INSTITUTE OF TECHNOLOGY \& SCIENCE, PILANI}
	\vspace*{0.2cm}
	\centerline{\bf Practice School Division}
	\vspace*{0.3cm}
	%\begin{center}
		\begin{tabular}{p{4.5cm} p{10.5cm}}
			\textbf{Station} & Software Tree, California, USA \\
			\textbf{Duration} & May 2026 -- July 2026 \\
			\textbf{Date of start} & 28 May 2026 \\
			\textbf{Date of submission} & 16 July 2026 \\
			\textbf{Title of the project} & Development of an Agentic Supply Chain Intelligence Platform using Gilhari, ORMCP, LangGraph, and Azure Cloud \\
			\textbf{ID No., Name and Discipline of the student} & 2024A7PS0085H, Sudhanva B H, B.E. Computer Science \\
			\textbf{Name and Designation of the expert} & Damodar Periwal, Founder/CEO, Software Tree LLC \\
			\textbf{Name of the PS faculty member} & Ajaya Kumar Pani \\
			\textbf{Key words} & Gilhari, JDX ORM, ORMCP, Model Context Protocol, LangGraph, LangChain, FastAPI, Flutter, Azure SQL Server, Azure Container Apps, Docker, Generative UI, Supply Chain, Agentic AI, REST API \\
			\textbf{Project area} & Agentic AI, Object-Relational Mapping, Cloud Computing, Full-Stack Development \\
			\textbf{Abstract} & This report covers the second and primary phase of PS-I work at Software Tree (May--July 2026). An initial Library Management project was completed first as practice with the Gilhari ORM toolchain. The main project then focused on building a full-stack Agentic Supply Chain Intelligence Platform. The system has five layers: a Microsoft Azure SQL Server database with a six-table normalized schema, a Dockerized Gilhari REST microservice deployed on Azure Container Apps, a Python ORMCP gateway exposing the database as typed MCP tools, a LangGraph agentic backend that runs a tool-calling loop with either OpenAI GPT-4o or Google Gemini, and a Flutter web dashboard with manual CRUD views and an AI Analytics chat panel that renders eleven different types of Generative UI responses. The complete platform is deployed on Microsoft Azure and the codebase is hosted on GitHub. \\
			& \\
			& \\
			& \\
			& \\
			& \\
			& \\
			\textbf{Signature of the student} & \hspace*{2cm}\textbf{Signature of the PS faculty member} \\
			& \\
			\textbf{Date} & \hspace*{2cm}\textbf{Date} \\
		\end{tabular}

	\newpage
	\centerline{\bf \Large Acknowledgements}
	\vspace*{0.5cm}

		\par\noindent I would like to thank Software Tree for the opportunity to work on Gilhari and agentic AI technologies during Practice School I.

		\par\noindent I am grateful to Mr.\ Damodar Periwal, Founder and CEO of Software Tree LLC, for his technical guidance and for the sessions on ORM concepts, the Gilhari framework, and the Model Context Protocol. His input was important in helping me understand the architecture and implement the system independently.

		\par\noindent I would also like to thank my mentor, Ms.\ Sruthi Gomathi Veeramani, for her support and coordination throughout the internship.

		\par\noindent I thank Prof.\ Ajaya Kumar Pani, Faculty-in-Charge and the Practice School Division, BITS Pilani, for organizing the programme. The internship gave me practical experience with cloud infrastructure, AI agent systems, and full-stack development that would be difficult to get from coursework alone.

		\vfill

		\par\noindent {\bf Signature of the student}

\tableofcontents
\listoffigures
\listoftables

% =========================================================
\chapter{Introduction}
% =========================================================

\section{Background}

PS-I at Software Tree was split into two phases. The first was a practice exercise: building a Library Management system using Gilhari from scratch. That involved defining a seven-class object model with nested BYVALUE relationships, writing the JDX ORM configuration, setting up PostgreSQL, containerizing the service with Docker, and building a Flutter dashboard. The purpose was to get comfortable with the Gilhari and JDX ORM toolchain before taking on the main assignment \cite{gilhari}.

Once that was done, the internship moved to the actual assigned work, which is what this report covers. The task was to build a production-level Agentic Supply Chain Intelligence Platform integrating cloud databases, Large Language Models, the Model Context Protocol, and a Generative UI system.

\section{Problem Statement}

Project 23, as described in the PS-I consent proforma, is titled \textit{``Developing an AI Agent with OpenAI Agents SDK, that uses ORMCP to access relational data on cloud. GPT, Azure, SQL Server.''}

The project description focuses on building an AI agent that uses OpenAI's developer tools and the Model Context Protocol to interact with a relational database through semantic, business-oriented interfaces. The official implementation tasks are:

\begin{itemize}[itemsep=0pt]
\item \textbf{Cloud Deployment}: Configuring, uploading, and running a Dockerized Gilhari microservice on the cloud using a free-tier developer license.
\item \textbf{Database Management}: Setting up and exchanging data with a Microsoft SQL Server database.
\item \textbf{Version Control and Documentation}: Uploading the codebase to GitHub with a comprehensive, professional README file.
\item \textbf{Technical Writing or Media Creation}: Publishing at least one technical article about the project on a platform such as Medium or LinkedIn, or creating a YouTube video detailing how ORMCP was integrated into the AI project, cross-linked with the GitHub repository.
\end{itemize}

The technology stack specified in the problem statement is: OpenAI GPT as the AI engine, the OpenAI Agents SDK as the agent framework, Microsoft Azure as the cloud platform, Microsoft SQL Server as the database, and ORMCP with the Gilhari microservice framework as the integration layer.

The implementation in this project covers all required tasks and goes further in several areas. In particular, the agent framework was implemented using LangChain and LangGraph rather than the bare OpenAI Agents SDK, which allowed the system to support both OpenAI GPT and Google Gemini through a single codebase. A full Flutter web admin dashboard was also built as an additional deliverable, providing both manual CRUD interfaces and an AI Analytics chat panel with eleven Generative UI response types. These additions are beyond the scope of the original problem statement but were built as part of the internship work.

\section{Objectives}

The technical objectives for this project were:
\begin{itemize}[itemsep=0pt]
\item Design and implement a six-table supply chain schema on Azure SQL Server.
\item Write the JDX ORM mapping and compile Java container classes for the Gilhari microservice.
\item Containerize Gilhari with Docker and deploy it to Azure Container Apps.
\item Set up the ORMCP server so the database is exposed as typed MCP tools.
\item Implement a LangGraph agent loop that calls those tools in response to natural language queries.
\item Design eleven Generative UI response types and implement them as Flutter widgets.
\item Build a Flutter admin dashboard supporting both manual CRUD and AI-driven analytics.
\item Deploy the full system to Azure: SQL Database, Container Apps, App Service, and Static Web Apps.
\item Validate everything end-to-end with integration tests and sample queries.
\end{itemize}

\section{Scope of the Report}

This report covers the primary project phase from late June 2026 to 16 July 2026. It describes all five layers of the system: the database, the Gilhari microservice, the ORMCP gateway, the FastAPI agent backend, and the Flutter dashboard. Both local development and Azure cloud deployment are covered. The cloud deployment, documented in \texttt{DEPLOYMENT.md}, is treated as complete and verified.


% =========================================================
\chapter{Literature Review and Technical Background}
% =========================================================

\section{Object-Relational Mapping and the Gilhari Framework}

Object-Relational Mapping (ORM) is a way to bridge object-oriented code and relational databases. Instead of writing SQL for every data operation, the developer writes a mapping that the ORM framework uses to handle persistence automatically. Gilhari, from Software Tree, is an ORM microservice built on the JDX engine \cite{gilhari}. Given compiled Java container classes and a JDX mapping file, it generates a full REST API (GET, POST, PUT, DELETE) for each mapped object with no data access code needed.

\section{JDX ORM Specification Language}

The JDX file (\texttt{.jdx}) is the configuration that drives Gilhari \cite{jdx}. It lists the database connection, the Java package, and for each class: the table name, attributes with Java types, primary key, and any relationships. JDX supports BYVALUE containment, where a parent object owns its child collection. In this project, a \texttt{PurchaseOrder} embeds an array of \texttt{PurchaseOrderItem} objects, and a \texttt{Supplier} contains its \texttt{InventoryItem} records.

\section{Model Context Protocol}

The Model Context Protocol (MCP) is an open specification that lets AI agents communicate with external tools and data sources in a standardized way \cite{mcp}. An MCP server registers a set of named, typed tools, each with a description and a JSON Schema for its parameters. A client (such as an agent framework) can list those tools, pick the right one, call it with structured arguments, and receive a structured result. This is how the AI agent in this project queries the database: through safe, schema-validated function calls rather than generated SQL.

\section{ORMCP: The Semantic Gateway}

ORMCP is a Python package from Software Tree that connects Gilhari to the MCP ecosystem \cite{ormcp}. When it starts, it reads the running Gilhari instance's object model and auto-generates an MCP tool for each entity: \texttt{query\_InventoryItems}, \texttt{insert\_StockTransaction}, \texttt{update\_PurchaseOrder}, \texttt{delete\_Supplier}, \texttt{execute\_raw\_sql}, and so on. These are served over standard I/O so any MCP client can discover and call them without writing any integration code.

\section{LangChain and LangGraph}

LangChain is a Python framework that provides a consistent interface for working with LLMs, tools, and prompt templates \cite{langchain}. It works with multiple providers (OpenAI, Google Gemini, Anthropic) through a common API, which is useful when you want to switch models without changing the agent logic. In this project it wraps the ORMCP tools as \texttt{StructuredTool} objects and enforces the response schema using \texttt{with\_structured\_output()}.

LangGraph sits on top of LangChain and handles stateful, multi-step agent loops \cite{langgraph}. It models the agent as a directed graph where each node does one thing (invoke the LLM, or run a tool call) and edges define what happens next. The full message history is kept in persistent state, so the agent can accumulate data over several tool calls before producing a final answer.

\section{FastAPI}

FastAPI is a Python web framework built on Starlette and Pydantic \cite{fastapi}. It handles async requests natively and generates OpenAPI documentation automatically. In this project it serves as the main entry point for the backend: handling chat requests, proxying Gilhari calls, verifying authentication, and streaming NDJSON responses to the frontend using \texttt{StreamingResponse}.

\section{Microsoft SQL Server and Azure SQL}

SQL Server is the database used throughout this project \cite{mssql}. Locally it runs as a Docker container (SQL Server 2022 Developer edition). In production, Azure SQL Database, a fully managed cloud version of SQL Server, is used instead. Gilhari communicates with it through the Microsoft JDBC driver (\texttt{mssql-jdbc-12.4.2.jre8.jar}), which is downloaded at Docker image build time.

\section{Microsoft Azure Cloud Infrastructure}

Azure is the cloud platform used for the production deployment \cite{azure}. Four services were used:
\begin{itemize}[itemsep=0pt]
\item \textbf{Azure SQL Database}: Managed cloud database for the six-table supply chain schema.
\item \textbf{Azure Container Apps}: Serverless container hosting for the Gilhari Docker image.
\item \textbf{Azure App Service}: Linux-based hosting for the Python FastAPI backend (Free F1 tier).
\item \textbf{Azure Static Web Apps}: Static hosting for the compiled Flutter web app, with GitHub Actions-based CI/CD.
\end{itemize}

\section{Flutter Web and Riverpod}

Flutter is Google's cross-platform UI framework that can compile to web JavaScript \cite{flutter}. The admin dashboard in this project is a Flutter web app. It uses Riverpod for state management, \texttt{fl\_chart} for bar, pie, and line charts, \texttt{flutter\_markdown} for rendering text responses, and \texttt{flutter\_svg} for the world map widget. The UI uses a dark theme with Google Fonts.

\section{Generative UI}

Generative UI is a pattern where the AI agent returns a structured JSON payload instead of free text, and the frontend renders an appropriate widget based on a type discriminator \cite{genui}. Rather than just describing data in words, the agent selects the right visualization (chart, Kanban board, KPI cards, form) and fills it with the actual query results. In this project the agent's output is validated against a Pydantic \texttt{AgentResponse} model with eleven possible response types, enforced by LangChain's \texttt{with\_structured\_output()} API.


% =========================================================
\chapter{Methodology}
% =========================================================

\section{System Architecture Design}

The platform has five layers with a clean separation of concerns, and two distinct ways for users to interact with it. Figure~\ref{fig:architecture} shows the overall structure.

\begin{figure}[ht]
\centering
\imgplaceholder{System Architecture Diagram}{Five-layer architecture with Pathway A (Manual CRUD: Flutter to Gilhari REST) and Pathway B (Agentic AI: Flutter to FastAPI to LangGraph to ORMCP to Gilhari to Azure SQL).}
\caption{Agentic Supply Chain Platform -- System Architecture}
\label{fig:architecture}
\end{figure}

\textbf{Pathway A (Manual CRUD)} is the straightforward one: the supply chain manager uses the data grid views in the Flutter dashboard to browse and edit records. These requests go directly to the Gilhari REST API through the FastAPI proxy, with no AI involved.

\textbf{Pathway B (Agentic Query)} is where the AI comes in. The user types a natural language question into the AI Analytics chat panel. The FastAPI backend passes it to the LangGraph agent, which calls ORMCP tools to query the database, accumulates the results, then asks the LLM to format a response. The LLM picks the right visualization type and fills in the data. The Flutter frontend receives this as a structured JSON object and renders the widget.

Table~\ref{tab:techstack} summarises which technology handles each layer.

\begin{table}[ht]
\centering
\begin{tabular}{|p{3cm}|p{4cm}|p{7cm}|}
\hline
\textbf{Layer} & \textbf{Technology} & \textbf{Role} \\
\hline
Frontend & Flutter Web + Riverpod & Admin dashboard: manual CRUD grids and AI Analytics chat panel with Generative UI \\
\hline
AI Orchestration & FastAPI + LangGraph + LangChain & Receives chat requests, runs the agent tool-calling loop, streams structured JSON responses \\
\hline
Semantic Gateway & ORMCP (\texttt{ormcp-server}) & Reads Gilhari's object model and exposes each entity as typed MCP tools \\
\hline
Data Abstraction & Gilhari (Docker) & Generates REST endpoints from JDX mappings; no hand-written data access code \\
\hline
Database & Microsoft SQL Server / Azure SQL & Stores the six-table supply chain schema \\
\hline
ORM Mapping & JDX + Java Container Classes & Defines the mapping from Java objects to SQL tables; compiled into the Gilhari container \\
\hline
\end{tabular}
\caption{Technology stack summary}
\label{tab:techstack}
\end{table}

\section{Database Schema Design}

The schema has six normalized tables covering the main entities in an inventory and procurement system. It was designed to have enough relational complexity to test both the manual CRUD layer and the agent's ability to navigate multi-table queries. Table~\ref{tab:schema} describes each table.

\begin{table}[ht]
\centering
\begin{tabular}{|p{3.5cm}|p{4.5cm}|p{6cm}|}
\hline
\textbf{Table} & \textbf{Purpose} & \textbf{Key Relationships} \\
\hline
\texttt{Suppliers} & Vendor profiles: company name, contact email, region & One-to-Many with \texttt{InventoryItems} \\
\hline
\texttt{ItemCategories} & Product category classifications & One-to-Many with \texttt{InventoryItems} \\
\hline
\texttt{InventoryItems} & Product catalog with stock quantities and unit prices & Many-to-One with \texttt{Suppliers} and \texttt{ItemCategories}; One-to-Many with \texttt{StockTransactions} \\
\hline
\texttt{PurchaseOrders} & Order headers with date, total cost, delivery status & One-to-Many with \texttt{PurchaseOrderItems} \\
\hline
\texttt{PurchaseOrderItems} & Junction table linking items to orders with negotiated prices & Many-to-One with \texttt{PurchaseOrders} and \texttt{InventoryItems} \\
\hline
\texttt{StockTransactions} & Audit log of stock changes (RESTOCK, SALE, ADJUSTMENT) & Many-to-One with \texttt{InventoryItems} \\
\hline
\end{tabular}
\caption{Six-table supply chain database schema}
\label{tab:schema}
\end{table}

\begin{figure}[ht]
\centering
\imgplaceholder{Entity-Relationship Diagram}{Six tables with their primary/foreign keys and the one-to-many relationships between them. Highlight the junction table (PurchaseOrderItems).}
\caption{Supply Chain Database Entity-Relationship Diagram}
\label{fig:dbschema}
\end{figure}

The schema was created by running \texttt{sql/init.sql} inside the SQL Server container. For the Azure deployment, a Python script (\texttt{init\_azure\_db.py}) was used instead, as it connects remotely via \texttt{pyodbc} and runs the same SQL file without needing local \texttt{sqlcmd} access. Test data was inserted using \texttt{seed.py}, which sends POST requests to the Gilhari REST API.

\section{JDX ORM Mapping and Java Container Classes}

Six Java container classes were created under the \texttt{com.supplychain.model} package, one for each table. Each extends \texttt{JDX\_JSONObject}. The \texttt{config/supply\_chain.jdx} file defines the mapping for all six, including the JDBC connection string, attribute types, primary keys, and BYVALUE relationships. A representative excerpt:

\begin{verbatim}
CLASS .Supplier TABLE Suppliers
    VIRTUAL_ATTRIB supplierID ATTRIB_TYPE int
    VIRTUAL_ATTRIB companyName ATTRIB_TYPE java.lang.String
    VIRTUAL_ATTRIB contactEmail ATTRIB_TYPE java.lang.String
    VIRTUAL_ATTRIB region ATTRIB_TYPE java.lang.String
    PRIMARY_KEY supplierID
    RELATIONSHIP inventoryItems
        REFERENCES ArrayInventoryItem BYVALUE WITH supplierID
;

CLASS .PurchaseOrder TABLE PurchaseOrders
    VIRTUAL_ATTRIB orderID ATTRIB_TYPE int
    VIRTUAL_ATTRIB orderDate ATTRIB_TYPE java.lang.String
    VIRTUAL_ATTRIB totalCost ATTRIB_TYPE double
    VIRTUAL_ATTRIB deliveryStatus ATTRIB_TYPE java.lang.String
    PRIMARY_KEY orderID
    RELATIONSHIP purchaseOrderItems
        REFERENCES ArrayPurchaseOrderItem BYVALUE WITH orderID
;
\end{verbatim}

The Java files were compiled with \texttt{compile.cmd} into the \texttt{bin/} directory. The \texttt{config/classnames\_map.json} file maps short class names to their fully qualified Java names, which Gilhari needs to resolve REST endpoints correctly.

\section{Docker Infrastructure}

The local setup uses Docker Compose with two services. The \texttt{sqlserver} service runs the official SQL Server 2022 image on port \texttt{1433}, with data persisted in a named volume (\texttt{sqlserver\_data}). The \texttt{gilhari-service} is a custom image extending \texttt{softwaretree/gilhari}. The \texttt{Dockerfile} downloads the Microsoft JDBC driver at build time using \texttt{wget}, then adds the compiled \texttt{bin/} classes, the \texttt{config/} directory, and \texttt{gilhari\_service.config}. Gilhari starts on internal port \texttt{8081}, mapped to port \texttt{80} on the host.

\section{ORMCP Semantic Gateway Setup}

The ORMCP server runs as a child process of the FastAPI app, managed through FastAPI's \texttt{lifespan} context. On startup it connects to Gilhari, reads the object model, and registers MCP tools for all six entities. The tool set covers query, insert, update, and delete for each entity, plus \texttt{execute\_raw\_sql} for joins and aggregations, and \texttt{getObjectModelSummary} for schema introspection. Connectivity was confirmed using \texttt{test\_mcp\_client.py}, which lists all tool schemas and writes them to \texttt{mcp\_tools\_output.log}.

\begin{figure}[ht]
\centering
\imgplaceholder{ORMCP Tool Discovery Output}{Output of test\_mcp\_client.py or contents of mcp\_tools\_output.log showing the auto-generated MCP tool names and their JSON Schema definitions.}
\caption{ORMCP Tool Discovery -- Available MCP Tools from the Supply Chain Model}
\label{fig:ormcp_tools}
\end{figure}

\section{LangGraph Agentic Backend}

The agent is in \texttt{backend/agent.py}, implemented as a LangGraph \texttt{StateGraph}. State holds the message history as a list of LangChain message objects (\texttt{HumanMessage}, \texttt{AIMessage}, \texttt{ToolMessage}). The graph has two nodes:

\textbf{Node 1: \texttt{call\_model}} calls the LLM with the current history and the full list of ORMCP tools wrapped as LangChain \texttt{StructuredTool} objects. The LLM provider is set by the \texttt{LLM\_PROVIDER} environment variable: \texttt{google} loads \texttt{ChatGoogleGenerativeAI} (Gemini Flash by default), and \texttt{openai} loads \texttt{ChatOpenAI} (GPT-4o-mini by default). Since both use LangChain's common interface, the agent logic is identical regardless of which provider is active. If the model returns tool calls, execution moves to \texttt{execute\_tools}. If it returns a plain message, the loop ends.

\textbf{Node 2: \texttt{execute\_tools}} runs each tool call through \texttt{MCPClient.call\_tool()}, appends the results as \texttt{ToolMessage} objects, and sends execution back to \texttt{call\_model}.

After the loop exits, the LLM is called once more with \texttt{with\_structured\_output(AgentResponse)} to produce the final Generative UI payload. This is then streamed to the Flutter frontend as NDJSON.

The system prompt tells the agent to call \texttt{getObjectModelSummary} first on complex queries, to prefer the named ORMCP tools for simple reads and \texttt{execute\_raw\_sql} for joins, and not to use raw SQL for any write operations.

\section{Generative UI Response System}

The \texttt{AgentResponse} model in \texttt{backend/models.py} is the contract between the agent and the Flutter frontend. It has a \texttt{response\_type} field and eleven possible payload classes, one for each widget type. Table~\ref{tab:genui} lists all of them.

\begin{table}[ht]
\centering
\begin{tabular}{|p{4.5cm}|p{9cm}|}
\hline
\textbf{Response Type} & \textbf{Rendered Widget in Flutter} \\
\hline
\texttt{text\_only} & Plain markdown text via \texttt{flutter\_markdown} \\
\hline
\texttt{table\_view} & Static scrollable \texttt{DataTable} with alternating row colors \\
\hline
\texttt{direct\_fetch\_table} & Live data table fetched directly from the Gilhari REST API \\
\hline
\texttt{metric\_kpi\_view} & Row of KPI cards with trend icons and color coding \\
\hline
\texttt{chart\_view} & Bar, pie, or line chart via \texttt{fl\_chart} \\
\hline
\texttt{timeline\_view} & Vertical event timeline with icon, timestamp, title, and subtitle \\
\hline
\texttt{kanban\_view} & Horizontally scrollable Kanban board with color-coded cards \\
\hline
\texttt{regional\_view} & SVG world map with regional data overlaid \\
\hline
\texttt{actionable\_form\_view} & Typed input form; submitting produces a \texttt{confirmation\_view} \\
\hline
\texttt{alert\_anomaly\_view} & Severity-badged alert cards for detected anomalies \\
\hline
\texttt{confirmation\_view} & Confirmation dialog that calls \texttt{POST /api/execute-tool} on confirm \\
\hline
\end{tabular}
\caption{Generative UI response types and their Flutter widget renderings}
\label{tab:genui}
\end{table}

Token usage, cost (calculated from \texttt{api\_costs.json}), and tool-call sequences are logged to daily JSON files in \texttt{backend/logs/} for each query.

\section{Flutter Admin Dashboard}

The dashboard is a single-page Flutter web app with a left sidebar. Riverpod handles state management. The app has two modes:

\textbf{Manual CRUD (Pathway A):} Each of the six entities has a dedicated view built on a reusable \texttt{AsyncDataGrid} widget with pagination, a toolbar for add/edit/delete, and a \texttt{BaseModal} for forms. The \texttt{PurchaseOrderForm} is the most involved since it manages nested \texttt{PurchaseOrderItems} in the same operation.

\textbf{AI Analytics (Pathway B):} The \texttt{agent\_chat\_view.dart} keeps a scrollable message list. When the user submits a query, the frontend streams the NDJSON response from \texttt{POST /api/agentic-chat}. Status lines like ``Calling tool: query\_InventoryItems...'' show up as progress indicators while the agent is working. The final line is parsed as an \texttt{AgentResponse} and handed to \texttt{DynamicMessageWidget}, which renders the appropriate widget.

\begin{figure}[ht]
\centering
\imgplaceholder{Flutter Admin Dashboard -- CRUD View}{Screenshot of the Flutter dashboard showing one of the six CRUD data grid views (e.g., Inventory Items or Purchase Orders) with the data table and toolbar visible.}
\caption{Flutter Admin Dashboard -- Manual CRUD View}
\label{fig:flutter_crud}
\end{figure}

\begin{figure}[ht]
\centering
\imgplaceholder{Flutter Admin Dashboard -- AI Analytics Chat}{Screenshot of the AI Analytics tab with a user query and a Generative UI widget as the agent response (e.g., a Kanban board or KPI cards view).}
\caption{Flutter Admin Dashboard -- AI Analytics Chat Panel with Generative UI Response}
\label{fig:flutter_ai}
\end{figure}

\section{Cloud Deployment on Azure}

The full system is deployed across four Azure services. The live endpoints are:
\begin{itemize}[itemsep=0pt]
\item \textbf{Flutter Frontend (Azure Static Web Apps):} \href{https://thankful-water-0f80cc410.7.azurestaticapps.net/}{thankful-water-0f80cc410.7.azurestaticapps.net}
\item \textbf{FastAPI Backend (Azure App Service):} \href{https://supply-chain-agentic-1234.azurewebsites.net/docs}{supply-chain-agentic-1234.azurewebsites.net}
\item \textbf{Gilhari Microservice (Azure Container Apps):} \texttt{gilhari-service-app.orangesand-d052af28.southindia.azurecontainerapps.io}
\end{itemize}

\textbf{Azure SQL Database.} The database was provisioned on the Basic tier. A Python script (\texttt{init\_azure\_db.py}) was used to initialize the schema remotely via \texttt{pyodbc}, running \texttt{sql/init.sql} against the Azure SQL endpoint. The firewall was configured to allow Azure services, and the JDX file was updated with the new server hostname.

\textbf{Gilhari on Azure Container Apps.} The Gilhari Docker image was built locally and pushed to an Azure Container Registry (ACR). A Container App was then created on the Consumption plan using this image, with HTTP ingress configured on port 80. The \texttt{supply\_chain.jdx} inside the image was already pointing to the Azure SQL hostname.

\textbf{FastAPI Backend on Azure App Service.} A Linux App Service (Free F1 tier) was created for the Python backend. Because Azure looks for \texttt{app.py} by default, the startup command was set explicitly to launch \texttt{main.py} with Uvicorn. \texttt{CORSMiddleware} was added to allow requests from the Flutter frontend's domain. The backend was deployed using zip deployment, and Azure's Oryx build engine installed \texttt{requirements.txt} automatically. Environment variables including the LLM API keys, the Gilhari base URL, and the application password were configured through the App Service settings panel.

\textbf{Flutter Dashboard on Azure Static Web Apps.} A Static Web App was created in the Azure Portal and connected to the GitHub repository. Azure's default Oryx builder does not support Flutter, so the auto-generated GitHub Actions workflow was modified to install Flutter and build the app using \texttt{flutter build web} before the deploy step, with the Azure build step instructed to skip its own build phase and use the pre-built output. Every push to the \texttt{main} branch now triggers a rebuild and redeploy automatically.

\textbf{End-to-End Verification.} The live Flutter dashboard was opened, authenticated, and used to run queries against the deployed stack. The agent responded correctly, pulling data from the Azure SQL Database through the Container Apps-hosted Gilhari service.

\begin{figure}[ht]
\centering
\imgplaceholder{Azure Portal Screenshot}{Azure Portal showing the deployed resource group with all four services: Azure SQL Database, Container App (Gilhari), App Service (FastAPI), and Static Web App (Flutter dashboard).}
\caption{Azure Portal -- Deployed Supply Chain Platform Resources}
\label{fig:azure_portal}
\end{figure}


% =========================================================
\chapter{Results and Discussions}
% =========================================================

\section{Completed Deliverables}

The following were completed during the primary phase of PS-I:
\begin{itemize}[itemsep=0pt]
\item A six-table Azure SQL Server schema with test data seeded via \texttt{seed.py}.
\item Six Java container classes compiled to \texttt{.class} files in \texttt{bin/}.
\item A JDX ORM specification (\texttt{supply\_chain.jdx}) with all table mappings and BYVALUE relationships.
\item A \texttt{Dockerfile} and \texttt{docker-compose.yml} for the local two-container setup.
\item ORMCP integration confirmed by \texttt{test\_mcp\_client.py}, exposing the database as typed MCP tools.
\item A LangGraph backend (\texttt{backend/agent.py}) that works with both OpenAI GPT-4o and Google Gemini Flash.
\item Eleven Generative UI response types as Pydantic models and Flutter widgets.
\item A Flutter admin dashboard with six CRUD views and an AI Analytics chat panel.
\item Full cloud deployment on Azure across four services, with live public URLs.
\item Per-query token and cost logging to daily JSON files in \texttt{backend/logs/}.
\item A GitHub repository with \texttt{README.md}, backend and frontend component READMEs, and a \texttt{DEPLOYMENT.md} with step-by-step deployment instructions.
\end{itemize}

\section{Gilhari REST API Results}

After starting Docker Compose and initializing the schema, Gilhari was running on port 80. The health endpoint confirmed it was ready:
\begin{verbatim}
GET http://localhost/gilhari/v1/health/check
-> {"status": "ok"}
\end{verbatim}

Gilhari created REST endpoints for all six entities from the JDX mapping alone. A GET request to \texttt{/gilhari/v1/Supplier} returned each supplier with its \texttt{inventoryItems} array embedded in the response. The same nested behavior worked for \texttt{/gilhari/v1/PurchaseOrder}, which returned each order with its \texttt{purchaseOrderItems}. No endpoint code was written manually.

\begin{figure}[ht]
\centering
\imgplaceholder{Gilhari REST API Response}{Screenshot of a curl or browser response from the Gilhari REST endpoint (e.g., GET /gilhari/v1/Supplier), showing the nested JSON with the BYVALUE inventoryItems array.}
\caption{Gilhari REST API -- Sample Nested JSON Response for Supplier Endpoint}
\label{fig:gilhari_api}
\end{figure}

\section{ORMCP Integration Results}

Running \texttt{python test\_mcp\_client.py} successfully connected to the ORMCP server and listed all generated tools. The \texttt{mcp\_tools\_output.log} file contains the full JSON Schema for each tool across all six entities: query, insert, update, and delete, plus \texttt{execute\_raw\_sql} and \texttt{getObjectModelSummary}. At FastAPI startup the ORMCP connection was also confirmed in the console:
\begin{verbatim}
Connecting to ORMCP Server via stdio...
Successfully connected to ORMCP Server.
Loaded environment. Provider: GOOGLE, Model: gemini-flash-lite-latest
\end{verbatim}

\section{AI Agent Demonstration}

The agent handled a range of natural language queries, picking tools and response types correctly. The prompts below, taken from the project's \texttt{PROMPTS.md} file, were used to test each of the eleven Generative UI response types. Screenshots of the rendered widgets should be placed in the figures that follow each group.

\subsection{Text and Table Responses}

\textit{``Give me a quick summary of how this supply chain database is structured.''} produces a plain \texttt{text\_only} markdown response, as the question is conversational and does not require a visualization. \textit{``Can you fetch a table of all purchase orders that have a status of 'Pending'?''} produces a \texttt{table\_view} response with a scrollable data table.

\begin{figure}[ht]
\centering
\imgplaceholder{Agent Response -- Text and Table Views}{Run the prompts: (1) ``Give me a quick summary of how this supply chain database is structured.'' and (2) ``Can you fetch a table of all purchase orders that have a status of Pending?'' Screenshot both responses side by side or stacked.}
\caption{Agent Generative UI -- Text and Table Responses}
\label{fig:genui_text_table}
\end{figure}

\subsection{KPI Metrics View}

\textit{``What are our key metrics? I want to see the total number of inventory items, total suppliers, and total stock transactions as KPI cards.''} produces a \texttt{metric\_kpi\_view} with color-coded KPI summary cards, each showing a value and trend direction.

\begin{figure}[ht]
\centering
\imgplaceholder{Agent Response -- KPI Metrics View}{Run the prompt: ``What are our key metrics? I want to see the total number of inventory items, total suppliers, and total stock transactions as KPI cards.'' Screenshot the rendered KPI cards.}
\caption{Agent Generative UI -- KPI Metrics View}
\label{fig:genui_kpi}
\end{figure}

\subsection{Timeline View}

\textit{``Can you give me a chronological timeline of the 5 most recent stock transactions?''} produces a \texttt{timeline\_view} with a vertical sequence of events, each showing the transaction type, item, quantity, and timestamp.

\begin{figure}[ht]
\centering
\imgplaceholder{Agent Response -- Timeline View}{Run the prompt: ``Can you give me a chronological timeline of the 5 most recent stock transactions?'' Screenshot the rendered vertical timeline widget.}
\caption{Agent Generative UI -- Timeline View}
\label{fig:genui_timeline}
\end{figure}

\subsection{Chart Views}

\textit{``Please analyze our inventory. Draw a bar chart showing the total quantity in stock for each item category.''} produces a \texttt{chart\_view} bar chart. \textit{``What is the breakdown of our inventory by category? Show me a pie chart.''} produces a pie chart of the same data.

\begin{figure}[ht]
\centering
\imgplaceholder{Agent Response -- Bar Chart and Pie Chart}{Run the prompts: (1) ``Draw a bar chart showing the total quantity in stock for each item category.'' and (2) ``What is the breakdown of our inventory by category? Show me a pie chart.'' Screenshot both charts.}
\caption{Agent Generative UI -- Bar Chart and Pie Chart Views}
\label{fig:genui_charts}
\end{figure}

\subsection{Regional Distribution View}

\textit{``Where are our suppliers located globally? Show me a regional distribution view.''} produces a \texttt{regional\_view} with the SVG world map and regional data overlaid.

\begin{figure}[ht]
\centering
\imgplaceholder{Agent Response -- Regional Distribution View}{Run the prompt: ``Where are our suppliers located globally? Show me a regional distribution view.'' Screenshot the rendered world map with supplier region overlays.}
\caption{Agent Generative UI -- Regional Distribution View}
\label{fig:genui_regional}
\end{figure}

\subsection{Kanban Board View}

\textit{``Show me the status of all active purchase orders as a Kanban board.''} produces a \texttt{kanban\_view} with columns for each delivery status and color-coded cards for each order.

\begin{figure}[ht]
\centering
\imgplaceholder{Agent Response -- Kanban Board View}{Run the prompt: ``Show me the status of all active purchase orders as a Kanban board.'' Screenshot the horizontally scrollable Kanban board with status columns and color-coded order cards.}
\caption{Agent Generative UI -- Kanban Board View}
\label{fig:genui_kanban}
\end{figure}

\subsection{Actionable Form and Confirmation Views}

\textit{``We are running low on RTX 4090 GPUs. Please provide a form to order 50 more from our best supplier.''} produces an \texttt{actionable\_form\_view} with pre-filled fields. Submitting the form produces a \texttt{confirmation\_view} that shows the pending mutation and requires the user to confirm before it is executed.

\begin{figure}[ht]
\centering
\imgplaceholder{Agent Response -- Form and Confirmation Views}{Run the prompt: ``We are running low on RTX 4090 GPUs. Please provide a form to order 50 more from our best supplier.'' Screenshot the form view, then submit it and screenshot the confirmation dialog.}
\caption{Agent Generative UI -- Actionable Form and Confirmation Views}
\label{fig:genui_form_confirm}
\end{figure}

\subsection{Alert and Anomaly View}

\textit{``Are there any risks in our supply chain right now? Show me an alert view of items with critical stock levels.''} produces an \texttt{alert\_anomaly\_view} with severity-badged cards listing the detected issues and suggested actions.

\begin{figure}[ht]
\centering
\imgplaceholder{Agent Response -- Alert and Anomaly View}{Run the prompt: ``Are there any risks in our supply chain right now? Show me an alert view of items with critical stock levels.'' Screenshot the rendered alert cards with severity badges.}
\caption{Agent Generative UI -- Alert and Anomaly View}
\label{fig:genui_alert}
\end{figure}

\section{Cloud Deployment Results}

All four Azure services came up and worked end-to-end:
\begin{itemize}[itemsep=0pt]
\item Azure SQL was reachable from both the local machine (for seeding) and from the Container Apps environment.
\item The Gilhari Container App responded to health checks and entity queries at its public HTTPS URL.
\item The App Service backend spawned the ORMCP child process correctly and forwarded agent queries through the Azure-hosted Gilhari instance.
\item The Static Web App loaded, authenticated, and rendered CRUD views and Generative UI responses against the live Azure stack.
\end{itemize}

\section{Discussion}

One thing that became clear during development is how much the Gilhari-ORMCP-LangGraph stack reduces the amount of glue code needed to connect an LLM to a database. The agent never writes SQL directly; it calls typed MCP functions with known schemas. This means there is no risk of the LLM generating malformed queries, and the system stays consistent even as the schema changes.

The dual-provider design worked well in practice. Switching from Gemini to OpenAI just meant changing \texttt{LLM\_PROVIDER} in the environment file. The agent, the tool-calling logic, and the structured output step were completely unchanged.

The Generative UI approach turned out to be much more useful than a plain chatbot. When the agent returned a Kanban board with color-coded purchase orders, the data was immediately understandable. A paragraph of text describing the same information would have taken more effort to read. Defining the response schema up front with Pydantic and enforcing it through \texttt{with\_structured\_output()} was essential: without it, the LLM's output format drifted noticeably between queries.

On the deployment side, the most non-obvious issue was getting Flutter to build correctly inside the Azure Static Web Apps GitHub Actions pipeline. Azure's default Oryx build system does not support Flutter, so the workflow had to be modified manually to install Flutter, run \texttt{flutter build web}, and then skip Azure's own build step. The FastAPI backend also required explicitly setting the startup command because Azure defaults to looking for \texttt{app.py} rather than \texttt{main.py}.

One issue that cost debugging time was the naming alignment required across the JDX class names, the \texttt{classnames\_map.json} file, the compiled \texttt{.class} files, and the ORMCP tool names. A single mismatch causes silent failures: Gilhari simply does not register the endpoint, or ORMCP exposes a tool with wrong field names. The \texttt{test\_mcp\_client.py} script was the most useful diagnostic tool for catching these problems early.


% =========================================================
\chapter{Conclusions and Future Work}
% =========================================================

\section{Conclusion}

The project produced a working five-layer Agentic Supply Chain Intelligence Platform, fully deployed on Azure. Starting from the Library Management practice project, the work moved through database design, JDX ORM configuration, Docker infrastructure, ORMCP integration, LangGraph agent implementation, Generative UI development, and cloud deployment. Each layer built on the previous one, and the final system handles both routine manual operations and complex natural language queries.

The combination of Gilhari's automatic REST generation and ORMCP's MCP tool auto-discovery meant very little custom data-access code was needed. LangGraph handled the multi-step agent loop reliably, and the Generative UI system gave the agent's answers a practical form that is immediately useful in a supply chain context. The cloud deployment on Azure showed that the architecture works beyond a local setup without major code changes.

\section{Future Work}

There are several directions this could go next:
\begin{itemize}[itemsep=0pt]
\item \textbf{WebSocket streaming}: Streaming intermediate tool-call results to the frontend in real time would reduce the waiting time for complex queries.
\item \textbf{Multi-agent setup}: Splitting the agent into specialized sub-agents (procurement, inventory, analytics) coordinated by a supervisor would handle more complex domain-specific queries better.
\item \textbf{Broader database support}: Testing the stack with PostgreSQL or MySQL would check whether the architecture is portable or tied to SQL Server.
\item \textbf{Proper authentication}: The current shared password is only suitable for a demo. OAuth 2.0 with user roles would be needed for a real deployment.
\item \textbf{Agent memory}: Storing conversation history in a vector database would allow the agent to carry context between sessions rather than starting fresh each time.
\end{itemize}


% =========================================================
\chapter{Recommendations}
% =========================================================

Based on the work and the issues encountered during it:

\begin{itemize}
\item \textbf{Get the JDX mapping right first}: The JDX file drives everything: Gilhari's REST surface, the ORMCP tool names, and ultimately what the LLM can call. Naming inconsistencies between the JDX class names, the \texttt{classnames\_map.json}, and the Java source files cause silent failures that are hard to trace. Getting this layer consistent before writing any higher-level code saves a lot of time.

\item \textbf{Run the ORMCP handshake test early}: \texttt{test\_mcp\_client.py} confirms that the ORMCP server is connected, the Gilhari introspection worked, and the expected tools are available. This takes a few seconds to run and should be the first integration test after starting the containers.

\item \textbf{Use Pydantic schemas for the Generative UI}: Trying to parse a free-form LLM response into a widget payload is unreliable. Define the response schemas in Pydantic from the start and use LangChain's \texttt{with\_structured\_output()} to enforce them. Without this, the output format varies enough to break the frontend renderer.

\item \textbf{Build for LLM provider flexibility from the start}: LangChain's model abstractions make it easy to support multiple providers without changing the agent code. This is worth setting up early, both for comparing model quality during development and for working around rate limits when one provider is unavailable.

\item \textbf{For Azure deployment, read the platform-specific docs}: Azure App Service and Static Web Apps both have defaults that do not match typical Python and Flutter project structures. The startup command for FastAPI needs to be set manually, and the Static Web Apps GitHub Actions workflow needs to be customized for Flutter. These are small issues but they take time if you encounter them without expecting them.
\end{itemize}

% Removed nocite
\bibliography{ref}
\addcontentsline{toc}{chapter}{Bibliography}

\begin{appendices}
\chapter{Project Repository Checklist}
The GitHub repository prepared for this project contains:
\begin{itemize}
\item \texttt{src/} directory with Java container class source files (\texttt{Supplier.java}, \texttt{InventoryItem.java}, \texttt{ItemCategory.java}, \texttt{PurchaseOrder.java}, \texttt{PurchaseOrderItem.java}, \texttt{StockTransaction.java}).
\item \texttt{bin/} directory with pre-compiled Java \texttt{.class} files.
\item \texttt{config/} directory with \texttt{supply\_chain.jdx} and \texttt{classnames\_map.json}.
\item \texttt{sql/} directory with \texttt{init.sql} for schema creation.
\item \texttt{Dockerfile} and \texttt{docker-compose.yml} for local development.
\item \texttt{gilhari\_service.config} for the Gilhari REST server.
\item \texttt{backend/} with the FastAPI backend: \texttt{main.py}, \texttt{agent.py}, \texttt{mcp\_client.py}, \texttt{models.py}, \texttt{requirements.txt}, \texttt{.env.example}.
\item \texttt{admin\_dashboard/} with the full Flutter web application.
\item \texttt{seed.py} and \texttt{seed\_lite.py} for database population.
\item \texttt{init\_azure\_db.py} for remote Azure SQL initialization via \texttt{pyodbc}.
\item \texttt{start\_ormcp.bat} for standalone ORMCP server launch.
\item \texttt{test\_mcp\_client.py} for ORMCP connectivity and tool schema validation.
\item \texttt{compile.cmd} and \texttt{build.cmd} build scripts.
\item \texttt{.github/workflows/} with the customized Azure Static Web Apps CI/CD pipeline.
\item \texttt{OVERVIEW.md}, \texttt{README.md}, \texttt{DEPLOYMENT.md}, and \texttt{PROMPTS.md}.
\end{itemize}

\chapter{Main Tools and Technologies Used}
\begin{itemize}
\item Gilhari SDK and JDX ORM Framework (Software Tree)
\item ORMCP Server Python Package (\texttt{ormcp-server}, Software Tree)
\item Model Context Protocol (MCP) Specification
\item LangChain and LangGraph (Python)
\item OpenAI API (GPT-4o, GPT-4o-mini)
\item Google Gemini API (Gemini Flash, Gemini Pro)
\item FastAPI and Uvicorn (Python)
\item Pydantic
\item Java (JDK 8+) for container class compilation
\item Microsoft SQL Server 2022 (Docker) and Azure SQL Database
\item Microsoft JDBC Driver for SQL Server
\item pyodbc (for remote Azure DB initialization)
\item Docker and Docker Compose
\item Microsoft Azure (Container Apps, App Service, Static Web Apps, Container Registry, SQL Database)
\item Azure CLI (\texttt{az})
\item Flutter SDK 3.x and Dart
\item Riverpod, \texttt{fl\_chart}, \texttt{flutter\_markdown}, \texttt{flutter\_svg}
\item Python 3.11+
\item Git, GitHub, GitHub Actions
\item Visual Studio Code
\end{itemize}

\end{appendices}

\end{document}
